\documentclass[11pt,letterpaper]{article}
\usepackage[margin=0.85in,headheight=16pt]{geometry}
\usepackage[T1]{fontenc}
\usepackage[utf8]{inputenc}
\usepackage{lmodern,microtype}
\usepackage{amsmath,amssymb,mathtools,booktabs,longtable,tabularx,array}
\usepackage[dvipsnames]{xcolor}
\usepackage{hyperref,fancyhdr,titlesec,enumitem,graphicx,tcolorbox}
\usepackage{lastpage}
\definecolor{ink}{HTML}{14263B}
\definecolor{sea}{HTML}{167C8A}
\definecolor{gold}{HTML}{C49747}
\definecolor{pale}{HTML}{EEF5F6}
\definecolor{graytext}{HTML}{516275}
\hypersetup{colorlinks=true,linkcolor=ink,urlcolor=sea,bookmarksopen=true,pdfauthor={Independent Educational Engineering Suite}}
\titleformat{\section}{\Large\bfseries\color{ink}}{\thesection}{0.7em}{}[\vspace{-0.3em}\color{gold}\titlerule]
\titleformat{\subsection}{\large\bfseries\color{sea}}{\thesubsection}{0.6em}{}
\setlist[itemize]{leftmargin=1.3em,itemsep=2pt,topsep=3pt}
\setlist[enumerate]{leftmargin=1.5em,itemsep=2pt,topsep=3pt}
\pagestyle{fancy}
\fancyhf{}
\fancyhead[L]{\small\sffamily\color{graytext}SECURITY ANALYSIS / ENGINEERING SUITE}
\fancyhead[R]{\small\sffamily\color{graytext}\textbf{\thepage}/\pageref{LastPage}}
\renewcommand{\headrulewidth}{0.3pt}
\renewcommand{\footrulewidth}{0pt}
\newtcolorbox{methodbox}[1][]{colback=pale,colframe=sea,boxrule=0.55pt,arc=1mm,title={\bfseries Method note},fonttitle=\sffamily,#1}
\newtcolorbox{sourcebox}[1][]{colback=white,colframe=gold,boxrule=0.55pt,arc=1mm,title={\bfseries Source boundary},fonttitle=\sffamily,#1}
\newcommand{\term}[1]{\textbf{\color{ink}#1}}
\newcommand{\src}{\textsf{\color{sea}[SOURCE]}}
\newcommand{\ext}{\textsf{\color{gold}[EXTENSION]}}
\newcommand{\gap}{\textsf{\color{red}[SOURCE GAP]}}
\newcommand{\answer}[1]{\par\smallskip\noindent\textbf{Check:} #1\par}
\newcommand{\cover}[3]{%
\begin{titlepage}\pagecolor{ink}\color{white}\vspace*{1.15in}
{\sffamily\large\color{gold}INDEPENDENT TECHNICAL DOCUMENTATION SUITE\par}\vspace{0.4in}
{\Huge\sffamily\bfseries #1\par}\vspace{0.15in}
{\Large\sffamily #2\par}\vspace{0.65in}
\color{gold}\rule{\textwidth}{2pt}\par\vspace{0.4in}
{\large #3\par}\vfill
{\small\sffamily Based on a critical reading of Graham and Dodd's \emph{Security Analysis}, Seventh Edition (2023). Original educational analysis; not affiliated with the publisher.\par}
\end{titlepage}\nopagecolor\color{black}}
\newcommand{\front}[3]{\cover{#1}{#2}{#3}\tableofcontents\clearpage
\begin{sourcebox}\src\ means a concept is attributable to the relevant chapter or seventh-edition contributor, summarized independently rather than reproduced. \ext\ means a modern method or engineering recommendation prepared for this suite, not an original Graham--Dodd rule. Historical examples and standards should not be treated as current market requirements. All numerical examples are synthetic, educational, and not securities recommendations.\end{sourcebox}}
\setlength{\parskip}{5pt}\setlength{\parindent}{0pt}
\begin{document}
\front{Security Analysis}{52-Chapter Study and Assessment Guide}{Original chapter-specific prompts, source-page navigation and examination rubric}
\section{How to use this study guide}
Each of the 52 entries contains a source-aligned orientation, one concrete evidence exercise, one challenge question, and an explicitly modern extension. These are \emph{independent study aids}, not replacements for reading the assigned source pages. Physical PDF page numbers are 1-indexed and may differ from the book's printed page labels. Historical supplementary chapters are present in the attached PDF except Chapter 36, which is listed but absent.
\begin{methodbox}For each chapter: identify the claim; reproduce a small calculation from a clean dataset; state a counterexample; and complete the independent modern extension. For grading, require a dated workbook, a one-page conclusion, and an explicit invalidation condition.\end{methodbox}

\section{Part 1: Survey and Approach}

\subsection{Chapter 1: The Scope and Limitations of Security Analysis. The Concept of Intrinsic Value}
\src\ \textit{PDF page 130; main text.}\par
\textbf{Read for:} Intrinsic value is an analytical estimate rather than an observable market quote; distinguish conclusions supported by facts from uncertainty that cannot be eliminated.\par
\textbf{Evidence task:} Document a price, a defensible value interval, and facts that could invalidate it.\par
\textbf{Challenge question:} Which assumptions would most severely change the investment conclusion?\par
\ext\ \textbf{Modern lens:} Express value as scenarios and a probability-free range before estimating probabilities.\par

\subsection{Chapter 2: Fundamental Elements in the Problem of Analysis. Quantitative and Qualitative Factors}
\src\ \textit{PDF page 144; main text.}\par
\textbf{Read for:} Separate quantitative evidence (assets, earnings, dividends, coverage) from qualitative judgments about prospects, management, and industry.\par
\textbf{Evidence task:} Construct two-column evidence inventory and rank each claim by reliability.\par
\textbf{Challenge question:} What cannot be inferred from an attractive accounting ratio?\par
\ext\ \textbf{Modern lens:} Create assumption provenance and a red-team checklist.\par

\subsection{Chapter 3: Sources of Information}
\src\ \textit{PDF page 158; main text.}\par
\textbf{Read for:} Financial statements, supplemental schedules, official disclosures, and independent corroboration differ in accessibility and trustworthiness.\par
\textbf{Evidence task:} Create a reproducible filing and source log with dates, units, and adjustments.\par
\textbf{Challenge question:} What would independently corroborate a management presentation?\par
\ext\ \textbf{Modern lens:} Use structured filings and version-controlled data transformations.\par

\subsection{Chapter 4: Distinctions Between Investment and Speculation}
\src\ \textit{PDF page 168; main text.}\par
\textbf{Read for:} Investment and speculation cannot be identified by asset labels alone; analyze the nature of the commitment and the evidentiary basis for safety and return.\par
\textbf{Evidence task:} Write separate investment and speculation theses for the same risky claim.\par
\textbf{Challenge question:} Can a conservative bond price still embody speculation?\par
\ext\ \textbf{Modern lens:} Apply mandate suitability and loss-tolerance gates.\par

\subsection{Chapter 5: Classification of Securities}
\src\ \textit{PDF page 180; main text.}\par
\textbf{Read for:} Economic claims, creditor seniority, participating rights, and speculative features are more analytically useful than a security's trade name.\par
\textbf{Evidence task:} Map seniority, cash-payment duties, voting rights, and conversion terms.\par
\textbf{Challenge question:} How could a preferred claim behave more like unsecured credit?\par
\ext\ \textbf{Modern lens:} Represent the security as a term-sheet object for valuation and stress testing.\par

\section{Part 2: Fixed-Value Investments}

\subsection{Chapter 6: The Selection of Fixed-Value Investments}
\src\ \textit{PDF page 218; main text.}\par
\textbf{Read for:} Select fixed-value securities by prioritizing safety of principal and income over prospective upside; distinguish contractual form from economic assurance.\par
\textbf{Evidence task:} Screen issuer cash generation, debt ranking, and obligations under adversity.\par
\textbf{Challenge question:} Why is a high coupon not evidence of investment quality?\par
\ext\ \textbf{Modern lens:} Build downside cash-flow and covenant headroom dashboards.\par

\subsection{Chapter 7: The Selection of Fixed-Value Investments: Second and Third Principles}
\src\ \textit{PDF page 231; main text.}\par
\textbf{Read for:} A bond judged safe only in favorable conditions is not conservatively selected; severity and history of business stress matter.\par
\textbf{Evidence task:} Reconstruct trough earnings and interest coverage through a recession.\par
\textbf{Challenge question:} What happens to principal protection if revenue falls 25 percent?\par
\ext\ \textbf{Modern lens:} Estimate cyclical recovery, refinancing capacity, and funding runways.\par

\subsection{Chapter 8: Specific Standards for Bond Investment}
\src\ \textit{PDF page 246; main text.}\par
\textbf{Read for:} Formal bond safety standards are exclusion filters, not promises of repayment; earnings coverage and business type require interpretation.\par
\textbf{Evidence task:} Apply an issuer-specific coverage policy and identify false positives.\par
\textbf{Challenge question:} Would the same threshold be suitable across utilities and cyclical issuers?\par
\ext\ \textbf{Modern lens:} Add forward covenant tests and normalized EBITDA-to-interest and FFO metrics.\par

\subsection{Chapter 9: Specific Standards for Bond Investment (Continued)}
\src\ \textit{PDF page 866; supplementary chapter.}\par
\textbf{Read for:} Security terms, maturity, collateral, and enforceability affect the investment claim beyond numerical coverage.\par
\textbf{Evidence task:} Summarize indenture provisions and collateral release conditions.\par
\textbf{Challenge question:} Which contract terms could negate attractive headline coverage?\par
\ext\ \textbf{Modern lens:} Run a modern indenture and collateral perfection review.\par

\subsection{Chapter 10: Specific Standards for Bond Investment (Continued)}
\src\ \textit{PDF page 257; main text.}\par
\textbf{Read for:} Property value, debt service capacity, and leverage offer distinct evidence of bond safety; property value alone is not a going-concern guarantee.\par
\textbf{Evidence task:} Reconcile asset encumbrances and debt outstanding.\par
\textbf{Challenge question:} When does collateral value fail to protect bondholders?\par
\ext\ \textbf{Modern lens:} Stress enterprise value and secured-debt recovery.\par

\subsection{Chapter 11: Specific Standards for Bond Investment (Continued)}
\src\ \textit{PDF page 883; supplementary chapter.}\par
\textbf{Read for:} The cushion of junior capitalization and observed market equity can complement, but not replace, independent credit judgment.\par
\textbf{Evidence task:} Build a capitalization stack from senior debt to common equity.\par
\textbf{Challenge question:} Why can a rapidly changing stock price mislead a credit analyst?\par
\ext\ \textbf{Modern lens:} Use stressed market-cap cushions and fundamental asset coverage.\par

\subsection{Chapter 12: Special Factors in the Analysis of Railroad and Public-Utility Bonds}
\src\ \textit{PDF page 895; supplementary chapter.}\par
\textbf{Read for:} Railroad and utility bonds demand industry-specific attention to assets, regulation, maintenance spending, and revenue resilience.\par
\textbf{Evidence task:} Compute maintenance burden and regulated cash coverage from historical reports.\par
\textbf{Challenge question:} What is lost when applying generic industrial metrics?\par
\ext\ \textbf{Modern lens:} Translate industry context to regulated infrastructure and project finance.\par

\subsection{Chapter 13: Other Special Factors in Bond Analysis}
\src\ \textit{PDF page 915; supplementary chapter.}\par
\textbf{Read for:} Consolidated holding-company results can obscure legal-entity cash restrictions and claims against subsidiaries.\par
\textbf{Evidence task:} Diagram legal entities and structural subordination.\par
\textbf{Challenge question:} Can parent debt be paid with trapped subsidiary cash?\par
\ext\ \textbf{Modern lens:} Use legal-entity liquidity and restricted-payment scenario analysis.\par

\subsection{Chapter 14: The Theory of Preferred Stocks}
\src\ \textit{PDF page 921; supplementary chapter.}\par
\textbf{Read for:} Preferred stock may combine capped participation with weak enforceability of distributions and principal repayment.\par
\textbf{Evidence task:} Compare preferred terms to a pari passu bond and common share.\par
\textbf{Challenge question:} Which promised return is discretionary rather than legally enforceable?\par
\ext\ \textbf{Modern lens:} Model cumulative dividends, call rights, and loss absorption.\par

\subsection{Chapter 15: Technique of Selecting Preferred Stocks for Investment}
\src\ \textit{PDF page 933; supplementary chapter.}\par
\textbf{Read for:} Require unusually robust issuer strength when a preferred security offers weaker creditor protections.\par
\textbf{Evidence task:} Perform coverage and senior-claim ranking review for preferred distributions.\par
\textbf{Challenge question:} Does bond-like yield compensate for preferred-specific subordination?\par
\ext\ \textbf{Modern lens:} Apply issuer and instrument-level cash-flow waterfall scenarios.\par

\subsection{Chapter 16: Income Bonds and Guaranteed Securities}
\src\ \textit{PDF page 945; supplementary chapter.}\par
\textbf{Read for:} Income bonds create contingent interest arrangements; guarantees require identification of the actual guarantor and its resources.\par
\textbf{Evidence task:} Evaluate payment trigger definitions, guarantee scope, and issuer coverage.\par
\textbf{Challenge question:} Is the guarantee independent of the obligor's distress?\par
\ext\ \textbf{Modern lens:} Model guarantor contagion and conditional coupons.\par

\subsection{Chapter 17: Guaranteed Securities (Continued)}
\src\ \textit{PDF page 267; main text.}\par
\textbf{Read for:} Mortgage and security guarantees vary materially in recourse, asset protection, and guarantor solvency.\par
\textbf{Evidence task:} Trace chain of obligation from mortgagor through guarantor.\par
\textbf{Challenge question:} What happens if guarantor and borrower fail simultaneously?\par
\ext\ \textbf{Modern lens:} Perform wrong-way risk and collateral valuation exercises.\par

\subsection{Chapter 18: Protective Covenants and Remedies of Senior Security Holders}
\src\ \textit{PDF page 281; main text.}\par
\textbf{Read for:} Protective covenants define creditor remedies, and rights are valuable only when enforceable and timely.\par
\textbf{Evidence task:} Mark key covenant language, notice periods, and acceleration triggers.\par
\textbf{Challenge question:} When could a covenant violation fail to prevent permanent loss?\par
\ext\ \textbf{Modern lens:} Prototype a monitoring calendar and covenant headroom alerts.\par

\subsection{Chapter 19: Protective Covenants (Continued)}
\src\ \textit{PDF page 294; main text.}\par
\textbf{Read for:} Negative pledges and additional-indebtedness restrictions affect changes in claim ranking after issuance.\par
\textbf{Evidence task:} Model incremental secured debt and dilution of existing debt claims.\par
\textbf{Challenge question:} What new liens could legally prime this security?\par
\ext\ \textbf{Modern lens:} Analyze baskets, restricted subsidiaries, and leakage capacity.\par

\subsection{Chapter 20: Preferred-Stock Protective Provisions. Maintenance of Junior Capital}
\src\ \textit{PDF page 958; supplementary chapter.}\par
\textbf{Read for:} Preferred protections can prohibit priming issues and preserve junior-capital backing, subject to approval provisions.\par
\textbf{Evidence task:} Extract issuance vetoes, coverage tests, and voting thresholds.\par
\textbf{Challenge question:} What dilution is possible without holder consent?\par
\ext\ \textbf{Modern lens:} Add pre- and post-recapitalization waterfall comparisons.\par

\subsection{Chapter 21: Supervision of Investment Holdings}
\src\ \textit{PDF page 305; main text.}\par
\textbf{Read for:} An investment is not automatically safe forever; quality must be supervised when facts or issuer conditions change.\par
\textbf{Evidence task:} Design a surveillance schedule and thesis-change protocol.\par
\textbf{Challenge question:} Which deterioration should trigger re-underwriting instead of passive holding?\par
\ext\ \textbf{Modern lens:} Automate watchlists with auditable exception records.\par

\section{Part 3: Senior Securities with Speculative Features}

\subsection{Chapter 22: Privileged Issues}
\src\ \textit{PDF page 332; main text.}\par
\textbf{Read for:} A senior security with conversion or other privileges mixes creditor protection with contingent equity participation.\par
\textbf{Evidence task:} Decompose straight claim and optionality from the term sheet.\par
\textbf{Challenge question:} At what price does the option account for most of market value?\par
\ext\ \textbf{Modern lens:} Perform fixed-income floor plus option sensitivity decomposition.\par

\subsection{Chapter 23: Technical Characteristics of Privileged Senior Securities}
\src\ \textit{PDF page 342; main text.}\par
\textbf{Read for:} Evaluate conversion privileges, call terms, redemption mechanics, and related technical conditions.\par
\textbf{Evidence task:} Tabulate bond floors, parity values, conversion ratios, and issuer call rights.\par
\textbf{Challenge question:} What changes when the issuer exercises a call?\par
\ext\ \textbf{Modern lens:} Build a payoff-diagram and term-event simulator.\par

\subsection{Chapter 24: Technical Aspects of Convertible Issues}
\src\ \textit{PDF page 356; main text.}\par
\textbf{Read for:} The effective conversion bargain can change as prices, anti-dilution adjustments, and redemption terms change.\par
\textbf{Evidence task:} Compute conversion parity and parity premium under three equity prices.\par
\textbf{Challenge question:} Can an apparently cheap convertible become expensive after a reset?\par
\ext\ \textbf{Modern lens:} Add state-dependent option valuation and stock-borrow costs.\par

\subsection{Chapter 25: Senior Securities with Warrants. Participating Issues. Switching and Hedging}
\src\ \textit{PDF page 971; supplementary chapter.}\par
\textbf{Read for:} Warrants, participation clauses, switching, and hedging alter upside participation and risk exposures.\par
\textbf{Evidence task:} Draw security cash flows, detachable warrant terms, and hedge legs.\par
\textbf{Challenge question:} Which risk remains after a nominal delta hedge?\par
\ext\ \textbf{Modern lens:} Study path dependence, dilution, basis, and transaction costs.\par

\subsection{Chapter 26: Senior Securities of Questionable Safety}
\src\ \textit{PDF page 366; main text.}\par
\textbf{Read for:} Questionable-safety senior securities require explicit attention to impaired principal and distressed outcomes.\par
\textbf{Evidence task:} Create recovery ranges with alternative going-concern and liquidation outcomes.\par
\textbf{Challenge question:} Can high expected yield coexist with unacceptable permanent-loss risk?\par
\ext\ \textbf{Modern lens:} Add legal priority, restructuring duration, and probability sensitivity.\par

\section{Part 4: Theory of Common-Stock Investment: The Dividend Factor}

\subsection{Chapter 27: The Theory of Common-Stock Investment}
\src\ \textit{PDF page 399; main text.}\par
\textbf{Read for:} Common-stock value rests on business ownership and future earning power; apparent cheapness alone is not a sound analytical conclusion.\par
\textbf{Evidence task:} Compare historical average earnings with market valuation and balance-sheet support.\par
\textbf{Challenge question:} How does uncertainty in earnings differ from fixed coupon risk?\par
\ext\ \textbf{Modern lens:} Construct owner earnings and business-franchise hypotheses.\par

\subsection{Chapter 28: Newer Canons of Common-Stock Investment}
\src\ \textit{PDF page 417; main text.}\par
\textbf{Read for:} Shifting common-stock investment doctrines can overpay for growth or undervalue stable earnings; neither historical averages nor projections are infallible.\par
\textbf{Evidence task:} Compare low-growth and high-growth valuation assumptions against actual evidence.\par
\textbf{Challenge question:} When is a high multiple justified, and when is it a fragile forecast?\par
\ext\ \textbf{Modern lens:} Use reverse DCF and competitive-advantage fade.\par

\subsection{Chapter 29: The Dividend Factor in Common-Stock Analysis}
\src\ \textit{PDF page 427; main text.}\par
\textbf{Read for:} Analyze the dividend record jointly with earnings and asset values; a payout is not itself proof of economic earning power.\par
\textbf{Evidence task:} Reconcile dividends to cash available for distribution over a cycle.\par
\textbf{Challenge question:} Could a high dividend be financed by leverage or asset sales?\par
\ext\ \textbf{Modern lens:} Bridge payout ratios with FCFE and buybacks.\par

\subsection{Chapter 30: Stock Dividends}
\src\ \textit{PDF page 983; supplementary chapter.}\par
\textbf{Read for:} Distinguish distributions of additional shares from transfers of economic cash; per-share figures must be adjusted.\par
\textbf{Evidence task:} Reconcile a stock dividend with unchanged aggregate equity value.\par
\textbf{Challenge question:} Did shareholders receive a genuine cash yield?\par
\ext\ \textbf{Modern lens:} Normalize total-return and split-adjusted per-share metrics.\par

\section{Part 5: Analysis of the Income Account: The Earnings Factor}

\subsection{Chapter 31: Analysis of the Income Account}
\src\ \textit{PDF page 478; main text.}\par
\textbf{Read for:} Recast the income account to understand recurring operations rather than taking reported net income at face value.\par
\textbf{Evidence task:} Develop a multi-period revenue-to-net-income earnings bridge.\par
\textbf{Challenge question:} Which line items reflect operating earning power?\par
\ext\ \textbf{Modern lens:} Use standardized operating profit reconciliation and audit trail.\par

\subsection{Chapter 32: Extraordinary Losses and Other Special Items in the Income Account}
\src\ \textit{PDF page 493; main text.}\par
\textbf{Read for:} Nonrecurring losses and special items can distort normal earning power, yet recurring write-offs must not be casually removed.\par
\textbf{Evidence task:} Classify restructuring, litigation, and inventory charges by recurrence.\par
\textbf{Challenge question:} When is an 'adjustment' actually a recurring cost of doing business?\par
\ext\ \textbf{Modern lens:} Reconcile GAAP to adjusted results and cash impacts.\par

\subsection{Chapter 33: Misleading Artifices in the Income Account. Earnings of Subsidiaries}
\src\ \textit{PDF page 504; main text.}\par
\textbf{Read for:} Subsidiary reporting and artificial accounting treatments may overstate parent earning power.\par
\textbf{Evidence task:} Tie consolidated earnings to cash remittable to shareholders.\par
\textbf{Challenge question:} Could earnings rise while distributable cash falls?\par
\ext\ \textbf{Modern lens:} Perform related-party, consolidation, and forensic-accounting tests.\par

\subsection{Chapter 34: The Relation of Depreciation and Similar Charges to Earning Power}
\src\ \textit{PDF page 522; main text.}\par
\textbf{Read for:} Depreciation and related noncash charges affect accounting earnings and economic maintenance requirements differently.\par
\textbf{Evidence task:} Compare reported D\&A with replacement-capital spending.\par
\textbf{Challenge question:} Why does adding back depreciation not create free cash flow?\par
\ext\ \textbf{Modern lens:} Model maintenance versus growth capital expenditures.\par

\subsection{Chapter 35: Public-Utility Depreciation Policies}
\src\ \textit{PDF page 994; supplementary chapter.}\par
\textbf{Read for:} Utility depreciation policies historically varied, complicating comparison of reported earnings and capital maintenance.\par
\textbf{Evidence task:} Recast two utility income accounts to a common depreciation assumption.\par
\textbf{Challenge question:} Does a longer asset life overstate distributable income?\par
\ext\ \textbf{Modern lens:} Compare regulated asset base, maintenance cycle, and normalized D\&A.\par

\subsection{Chapter 36: Amortization Charges from the Investor's Standpoint}
\gap\ \textbf{Not present in the supplied PDF.} The title is verified in the contents; the orientation below is a limitation notice, not a chapter summary.\par
\textbf{Read for:} The supplied PDF's table of contents identifies this chapter, but its chapter text is absent; no detailed historical summary can be validated from the supplied file.\par
\textbf{Evidence task:} Locate the publisher's original supplementary chapter before performing a source-grounded close reading.\par
\textbf{Challenge question:} How does amortization of an asset differ from amortization of bond premium?\par
\ext\ \textbf{Modern lens:} Independent extension only: build schedules for amortized debt premiums, intangible charges, and cash-flow effects.\par

\subsection{Chapter 37: Significance of the Earnings Record}
\src\ \textit{PDF page 541; main text.}\par
\textbf{Read for:} A long earnings record can reveal persistence, volatility, and shifts that a single annual figure hides.\par
\textbf{Evidence task:} Compute medians, trough years, and cycle-normalized margins.\par
\textbf{Challenge question:} How many years are truly representative of recurring performance?\par
\ext\ \textbf{Modern lens:} Apply regime-change and rolling-window sensitivity tests.\par

\subsection{Chapter 38: Specific Reasons for Questioning or Rejecting the Past Record}
\src\ \textit{PDF page 556; main text.}\par
\textbf{Read for:} Past earnings may be nonrepresentative after resource depletion, business change, price shifts, or impaired competitive position.\par
\textbf{Evidence task:} Document four reasons to reject historical extrapolation for a hypothetical issuer.\par
\textbf{Challenge question:} Which causal drivers can invalidate a five-year average?\par
\ext\ \textbf{Modern lens:} Use unit economics and change-point analysis.\par

\subsection{Chapter 39: Price-Earnings Ratios for Common Stocks. Adjustments for Changes in Capitalization}
\src\ \textit{PDF page 565; main text.}\par
\textbf{Read for:} A capitalization multiple is meaningful only alongside dependable earning power and changes in shares or claims.\par
\textbf{Evidence task:} Reconcile per-share earnings after a capital raise and estimate justified P/E intervals.\par
\textbf{Challenge question:} What happens to the P/E when one-time gains are removed?\par
\ext\ \textbf{Modern lens:} Compare forward normalized earnings, cash yield, and ROIC reinvestment.\par

\subsection{Chapter 40: Capitalization Structure}
\src\ \textit{PDF page 576; main text.}\par
\textbf{Read for:} Priority claims, leverage, and different security classes change the allocation of total enterprise earnings to common shareholders.\par
\textbf{Evidence task:} Prepare a three-company capital structure sensitivity comparison.\par
\textbf{Challenge question:} Can identical enterprise earnings produce radically different EPS risk?\par
\ext\ \textbf{Modern lens:} Introduce enterprise-to-equity bridges and downside leverage effects.\par

\subsection{Chapter 41: Low-Priced Common Stocks. Analysis of the Source of Income}
\src\ \textit{PDF page 589; main text.}\par
\textbf{Read for:} Low nominal share price carries no intrinsic valuation information; examine how profits arise and claims are structured.\par
\textbf{Evidence task:} Compare low-priced and high-priced shares using per-share and enterprise metrics.\par
\textbf{Challenge question:} Is a one-dollar share less expensive than a hundred-dollar share?\par
\ext\ \textbf{Modern lens:} Normalize share count, enterprise value, and source-of-profit quality.\par

\section{Part 6: Balance-Sheet Analysis: Implications of Asset Values}

\subsection{Chapter 42: Balance-Sheet Analysis. Significance of Book Value}
\src\ \textit{PDF page 614; main text.}\par
\textbf{Read for:} Book value provides a starting balance-sheet anchor, subject to asset quality, accounting measurement, and liabilities.\par
\textbf{Evidence task:} Recalculate tangible book value after asset write-downs.\par
\textbf{Challenge question:} Why is a book-value discount sometimes justified?\par
\ext\ \textbf{Modern lens:} Incorporate intangible-heavy business models and off-balance commitments.\par

\subsection{Chapter 43: Significance of the Current-Asset Value}
\src\ \textit{PDF page 625; main text.}\par
\textbf{Read for:} Net current asset value focuses on liquid assets after deduction of liabilities and can indicate asset-backed undervaluation.\par
\textbf{Evidence task:} Compute NCAV, cash conversion quality, and available margin of safety.\par
\textbf{Challenge question:} Which current assets may not realize their stated accounting amount?\par
\ext\ \textbf{Modern lens:} Stress receivables, inventories, and lease/contingent liabilities.\par

\subsection{Chapter 44: Implications of Liquidating Value. Stockholder-Management Relationships}
\src\ \textit{PDF page 641; main text.}\par
\textbf{Read for:} Liquidation-value discounts depend on the ability to realize assets and on stockholder--management incentives.\par
\textbf{Evidence task:} Build a liquidation waterfall including costs and preferred claims.\par
\textbf{Challenge question:} Who can compel a sale, and what prevents value realization?\par
\ext\ \textbf{Modern lens:} Model catalysts, liquidation costs, and governance scenarios.\par

\subsection{Chapter 45: Balance-Sheet Analysis (Concluded)}
\src\ \textit{PDF page 657; main text.}\par
\textbf{Read for:} Balance-sheet analysis also reveals overvaluation, hidden weaknesses, and asset encumbrances, not only undervalued liquid assets.\par
\textbf{Evidence task:} Review liabilities, pledged assets, receivable quality, and capitalization ratios.\par
\textbf{Challenge question:} Where might apparently strong book equity fail to protect holders?\par
\ext\ \textbf{Modern lens:} Add asset impairments and contingent claims to downside cases.\par

\section{Part 7: Additional Aspects of Security Analysis: Price---Value Discrepancies}

\subsection{Chapter 46: Stock-Option Warrants}
\src\ \textit{PDF page 1009; supplementary chapter.}\par
\textbf{Read for:} Warrants provide long-dated equity option rights whose valuation depends on exercise price, time, and dilution.\par
\textbf{Evidence task:} Reconcile number of issued warrants with fully diluted share count.\par
\textbf{Challenge question:} Does apparent upside survive exercise dilution and expiry?\par
\ext\ \textbf{Modern lens:} Build a simplified option and dilution sensitivity model.\par

\subsection{Chapter 47: Cost of Financing and Management}
\src\ \textit{PDF page 698; main text.}\par
\textbf{Read for:} Promoter compensation, financing fees, and managerial arrangements can divert value before it reaches investors.\par
\textbf{Evidence task:} Trace fees, expenses, related-party transactions, and ownership layering.\par
\textbf{Challenge question:} Which claims capture value ahead of outside shareholders?\par
\ext\ \textbf{Modern lens:} Analyze total economic fee drag and minority-holder protections.\par

\subsection{Chapter 48: Some Aspects of Corporate Pyramiding}
\src\ \textit{PDF page 709; main text.}\par
\textbf{Read for:} Corporate pyramiding allows concentrated control over assets with comparatively small direct economic ownership.\par
\textbf{Evidence task:} Diagram voting control, economic interest, and intercompany leverage.\par
\textbf{Challenge question:} What rights does a minority investor really possess?\par
\ext\ \textbf{Modern lens:} Stress look-through leverage and governance discount.\par

\subsection{Chapter 49: Comparative Analysis of Companies in the Same Field}
\src\ \textit{PDF page 719; main text.}\par
\textbf{Read for:} Industry comparisons are useful only after aligning accounting, operations, capital intensity, and business mix.\par
\textbf{Evidence task:} Construct an adjustment ledger for three peer companies.\par
\textbf{Challenge question:} Does the cheapest trading multiple imply the best bargain?\par
\ext\ \textbf{Modern lens:} Use ROIC-adjusted peer sets and scenario-conditioned multiples.\par

\subsection{Chapter 50: Discrepancies Between Price and Value}
\src\ \textit{PDF page 734; main text.}\par
\textbf{Read for:} Price can diverge from carefully estimated value; investigate catalysts, neglected securities, and permanent impairments.\par
\textbf{Evidence task:} Document three alternative explanations for a 40 percent apparent discount.\par
\textbf{Challenge question:} Which explanation would make the discount justified rather than exploitable?\par
\ext\ \textbf{Modern lens:} Maintain a value-gap thesis with identified resolution pathways.\par

\subsection{Chapter 51: Discrepancies Between Price and Value (Continued)}
\src\ \textit{PDF page 753; main text.}\par
\textbf{Read for:} Seasoned versus lesser-known issues may trade differently even with similar fundamentals; neglected-credit anomalies require discrimination.\par
\textbf{Evidence task:} Compare bond terms, issuer operating evidence, and liquidity across two credits.\par
\textbf{Challenge question:} Is the spread premium compensation for default, illiquidity, or unfamiliarity?\par
\ext\ \textbf{Modern lens:} Decompose credit spreads into expected loss and risk/liquidity premia.\par

\subsection{Chapter 52: Market Analysis and Security Analysis}
\src\ \textit{PDF page 762; main text.}\par
\textbf{Read for:} Price-chart forecasting and fundamental security analysis ask different questions; technical signals are not substitutes for business value evidence.\par
\textbf{Evidence task:} Separate momentum observations from an investment-underwriting memorandum.\par
\textbf{Challenge question:} When does price action justify updating intrinsic-value assumptions?\par
\ext\ \textbf{Modern lens:} Introduce decision rules for information updates without equating price to value.\par

\section{Part VIII: contemporary contributor study units}
\subsection{International Investing (PDF page 778)}
\src\ Assess whether familiar valuation principles transfer to unfamiliar legal, governance, monetary and reporting environments. \textbf{Exercise:} create a foreign-security checklist with currency, transfer, ownership, regulatory, disclosure, and minority-holder risks. \ext\ Compare local-currency and investor-currency scenario outcomes.
\subsection{At the Intersection of Public and Private Investing: A Case Study (PDF page 798)}
\src\ Examine patient capital, private-market valuation, and the danger of jeopardizing financial survival. \textbf{Exercise:} construct a public/private ownership choice memo. \ext\ Add illiquidity lockup and transaction-cost sensitivity.
\subsection{Endowment Management Principles and Practice (PDF page 810)}
\src\ Investigate long-horizon allocation subject to spending requirements and institutional constraints. \textbf{Exercise:} reconcile expected liquidity demands with lockup commitments. \ext\ Stress a multi-year drawdown while annual spending continues.

\section{Cross-chapter synthesis and oral examination}
\begin{enumerate}
\item Defend a discounted bond after showing worst-year coverage, legal claim rank, and required refinancing.
\item Reconcile an apparently low P/E with normalized profits, lease/debt obligations, and dilution.
\item Explain when a stock below NCAV can nevertheless destroy capital.
\item Evaluate two securities of the same issuer with different guarantees and conversion rights.
\item Distinguish an attractive security from a suitable addition to an already concentrated portfolio.
\item Recalculate the investment case when the original high-margin year is found to contain one-time proceeds.
\end{enumerate}
\section{Evidence rubric}
\begin{center}\begin{tabularx}{\textwidth}{p{0.26\textwidth}p{0.16\textwidth}X}\toprule Dimension & Weight & Passing evidence \\
\midrule Source accuracy & 25\% & Exact chapter, instrument, and source page or filing identified \\
Calculation & 25\% & Reproducible figures, definitions, unit/date checks \\
Adversarial reasoning & 25\% & Explicit countercase, failure state, and remedy \\
Communication & 15\% & Concise decision, uncertainty disclosed \\
Modern translation & 10\% & Extension labeled and not attributed to Graham--Dodd \\
\bottomrule\end{tabularx}\end{center}

\end{document}
